
\image{frama-c}


\levelThreeTitle{Frama-C? WP?}



Frama-C (FRAmework for Modular Analysis of C code) is a platform
dedicated to the analysis of C programs created by the CEA LIST and
Inria. It is based on a modular architecture allowing you to use different
plugins. The default plugins comprise different static analyses (that
do not execute source code), dynamic analyses (that requires code
execution), or a combination of both. These plugins can collaborate together
or not, either by communicating directly or by using the specification
language provided by Frama-C.


\image{gallery}



This specification language is called ACSL (``Axel'') for
ANSI C Specification Language and allows us to express the
properties we want to verify about our programs. These properties are
written using code annotations in comment sections. If one has
already used Doxygen, it is quite similar, except that we write logic
formulas and not text. During this tutorial, we will extensively write
ACSL code, so let us just skip this for now.



The analysis we will use in this tutorial is provided by the WP plugin
(for Weakest Precondition), a deductive verification plugin. It implements
the technique we mentioned earlier: from
ACSL annotations and the source code, the plugin generates what we call
verification conditions, that are logic formulas that must be verified to be
satisfiable or not. This verification can be performed manually or
automatically; here we use automatic tools.



We will use an SMT solver
(\externalLink{satisfiability modulo theory}{https://fr.wikipedia.org/wiki/Satisfiability\_modulo\_theories}, we will not detail how it works). This solver is
\externalLink{Alt-Ergo}{http://alt-ergo.ocamlpro.fr/}, that was initially developed
by the Laboratoire de Recherche en Informatique d'Orsay, and is today
maintained by OCamlPro.



\levelThreeTitle{Installation}



Frama-C is a software developed on Linux and macOS. It is thus better supported
on these operating systems. Windows is currently supported through Windows
Subsystem for Linux (WSL), Frama-C should behave the same way as on Linux, but
this usage is less tested than pure Linux or macOS.


Different proof tools are available to realize proofs with Frama-C. Most of the
time, Frama-C and WP are used \textit{only} with automatic solvers, even if one
can also use interactive provers like Rocq. Using Rocq with Frama-C requires
installing Frama-C with Opam (because of how Why3 and Rocq interact).


Some proofs of this tutorial are finalized using Rocq but without detailing
them since it is not the purpose of this tutorial. It is not \textit{necessary}
to install Rocq to follow this tutorial, a few proofs will not be \textit{entirely}
finished, but this is not a problem.


\levelFourTitle{Installation packages}


Frama-C provides installation packages for Linux (and WSL) and macOS. These
packages allow installing Frama-C and its graphical user interface.


\levelFiveTitle{Linux \& WSL}


From \externalLink{the download page of Frama-C}{https://frama-c.com/html/get-frama-c.html},
get the installation package. From a terminal, make it executable and run
it :


\begin{CodeBlock}{shell}
$ chmod +x frama-c-linux-x86-64-33.0-Arsenic.run
$ sudo ./frama-c-linux-x86-64-33.0-Arsenic.run
\end{CodeBlock}


On Ubuntu, it is highly recommenced to execute the installation with \CodeInline{sudo}
so that the installation script can add the necessary permissions in AppArmor for
the graphical user interface. Else, running the interface will either require to
build an AppArmor profile, or disable the sandbox by always adding the option
\CodeInline{--no-sandbox}.


The same procedure can be applied in WSL, notice that depending on the WSL
distribution, it can miss basic libraries for graphical usage. For example,
a fresh WSL Debian may need different additional libraries to run the graphical
user interface. The default WSL Ubuntu embeds all necessary packages.


After installing Frama-C, we also need an automatic proof tool, see
Section~\ref{l5:program-proof-and-our-tool-frama-c-install-packages-solver}.


\levelFiveTitle{macOS}


For macOS, the installation package is split into two components, the core
component and the graphical user interface that must be installed separately.


From \externalLink{the download page of Frama-C}{https://frama-c.com/html/get-frama-c.html},
get:
\begin{itemize}
  \item the \CodeInline{.pkg} that corresponds to the target architecture,
  \item the graphical user interface.
\end{itemize}


After opening the \CodeInline{.pkg} (either via the graphical user interface, or
via the command line using \CodeInline{open}), just follow the installation
procedure. Similarly, the DMG application can be installed via the graphical
user interface, typically in \CodeInline{/Applications}.


After installing Frama-C, we also need an automatic proof tool, see
Section~\ref{l5:program-proof-and-our-tool-frama-c-install-packages-solver}.


\levelFiveTitle{Alt-Ergo}
\label{l5:program-proof-and-our-tool-frama-c-install-packages-solver}


In this tutorial, a single tool is absolutely necessary : Alt\nobreakdash-Ergo.
In order to install it, go on its \externalLink{release's page}{https://github.com/OCamlPro/alt-ergo/releases}.


The simplest way to install it is to take a pre-compiled binary that corresponds
to the target system, to rename it \CodeInline{alt-ergo} and to copy it in a
location that is part of the \CodeInline{PATH} (for example \CodeInline{/usr/bin}) ;
\begin{CodeBlock}{shell}
sudo mv alt-ergo-v2.6.4-x86_64-linux-musl /usr/bin/alt-ergo
\end{CodeBlock}


Go to the section ``Verify installation'' to perform some tests on
your installation.


\levelFourTitle{Opam \& NPM}


A second option is to use Opam, a package manager for Ocaml libraries
and applications, and NVM/NPM to compile the graphical user interface.


First, Opam must be installed (see
\externalLink{its documentation}{https://opam.ocaml.org/doc/Install.html}). Then,
one can run:


\begin{CodeBlock}{bash}
opam install frama-c
\end{CodeBlock}


This command installs Frama-C, and a compilation and installation script for the
Frama-C GUI. To run this compilation and installation, \CodeInline{nvm},
\CodeInline{npm} and \CodeInline{yarn} must be available, see
\externalLink{the documentation of NVM}{https://github.com/nvm-sh/nvm\#installing-and-updating}
to install NVM, then we can configure it to compile and install the GUI by just
running it a first time:


\begin{CodeBlock}{shell}
$ nvm install 24       # we need node 24
$ nvm use 24           # we use it
$ npm install -g yarn  # we need yarn
$ frama-c-gui          # compile and install the GUI
\end{CodeBlock}


On Ubuntu, AppArmor has a strict configuration, to run the GUI without additional
configuration, one has to disable the sandbox when running the interface using
the option \CodeInline{--no-sandbox}.


Go to the section ``Verify installation'' to perform some tests on
your installation.


\levelThreeTitle{Verify installation}


In order to verify that the installation has been correctly performed,
we will use the following code:



\CodeBlockInput{c}{verify.c}


Place this code in a file named \CodeInline{main.c}.
(You can find the above code at the GitHub repository under the
\CodeInline{code/program-proof-and-our-tool/frama-c/} subdirectory
if you would like to copy-paste it.)
Then, from a terminal, in the folder where the file has been created, we
start Frama-C with the following command line:



\begin{CodeBlock}{bash}
frama-c-gui -wp -wp-rte main.c
\end{CodeBlock}



The following window should appear:



\image{verif_install-1}


After selecting the WP view (which is dedicated to the WP plug-in), one can
select the ``Files'' side panel that lists the files and functions of the
program under analysis. We can see the (slightly modified) content of the
functions, and some green bullets on different lines as illustrated here:



\image{verif_install-2}


If this does not succeed, check that the graphical user interface is installed,
and the Alt-Ergo solver is available. If everything seems to be OK, forums can
help (for example: \externalLink{StackOverflow - Frama-C}{https://stackoverflow.com/questions/tagged/frama-c})



\begin{Warning}
  The graphical user interface of Frama-C does not allow editing source code.
\end{Warning}


\levelThreeTitle{(Bonus) Some more provers}


This part is optional, nothing in this section should be particularly
useful \emph{in the tutorial}. However, when we start to be interested
in proving more complex programs, it is often possible to reach the
limits of Alt-Ergo, which is available by default, and we would thus need
some other provers. For basic properties, almost all solvers have the
same capabilities, for more complex ones, each solver has its specialized
domains.


\levelFourTitle{Why3}


\externalLink{Why3}{https://why3.org/} is a deductive proof platform developed
by the Laboratoire Méthodes Formelles (LMF) in Orsay. It provides a programming
language and a specification language, and allows interacting with a wide
variety of automatic and interactive provers. This point is the one that
interests us here. WP uses Why3 as a backend to talk with external provers.


On their website, we can find the list of
\externalLink{supported provers}{https://www.why3.org/\#provers}.
We recommend installing
\externalLink{Z3}{https://github.com/Z3Prover/z3/wiki} which is
developed by Microsoft Research, and
\externalLink{CVC5}{https://cvc5.github.io/} which is developed by many
research teams (New York University, University of Iowa, Google, CEA
List). Those two provers are very efficient and somewhat complementary.


New provers can be installed any time after the installation of Frama-C. They
are then available in the WP sidebar in the ``Provers Configuration'' section:


\image{detect-and-select}


\levelFourTitle{Rocq}


Rocq, which is developed by Inria, is a proof assistant. Basically, we
write the proofs ourselves in a dedicated language and the assistant
verifies (using typing) that the proof is actually a valid proof.



Why would we need such a tool? Sometimes, the properties we want to
prove can be too complex to be solved automatically by SMT solvers,
typically when they require careful inductive reasoning with precise
choices at each step. In this situation, WP allows us to generate
verification conditions translated in Rocq language, and to write the
proof ourselves.


To learn Rocq, we would recommend
\externalLink{this tutorial}{http://www.cis.upenn.edu/~bcpierce/sf/current/index.html}.


If one needs more information about Rocq and its installation, this page
can help: \externalLink{The Rocq Proof Assistant}{https://rocq-prover.org/}.
